\section*{Capítulo \ref{ch:b3:statistical-thermo-applied} --- Termodinâmica estatística aplicada}

\begin{solution}{exo:b3:statistical-thermo-applied:1}
$\frac32R + R + (3 \times 3 - 5)R = 6.5R = \qty{54.0}{J.K^{-1}.mol^{-1}}$. À temperatura ambiente, a translação
e a rotação são clássicas, os dois estiramentos estão congelados e a deformação
angular duplamente degenerada, a vibração mais baixa, está parcialmente excitada.
\end{solution}

\begin{solution}{exo:b3:statistical-thermo-applied:2}
$x = 1$: $\eu/(\eu - 1)^2 = 0.921$. $x = 10$: $100\eu^{10}/(\eu^{10} - 1)^2 \approx 100\eu^{-10} = \num{4.5e-3}$.
\end{solution}

\begin{solution}{exo:b3:statistical-thermo-applied:3}
\qty{20}{K}: $x_{\mathrm{para}} = 1/(1 + 3 \times 3\eu^{-2 \times 85.35/20}) = 1/(1 + 9 \times \num{1.96e-4}) = 0.998$.
\qty{77}{K}: soma par $1 + 5\eu^{-6.651} = 1.0065$; soma ímpar $3(3\eu^{-2.217} + 7\eu^{-13.30}) = 0.9806$;
$x_{\mathrm{para}} = 1.0065/1.9871 = 0.51$.
\end{solution}

\begin{solution}{exo:b3:statistical-thermo-applied:4}
Para $I = 1$, há $3 \times 2 = 6$ estados de spin simétricos e $3 \times 1 = 3$ antissimétricos;
os dêuterons são bósons, logo a função de onda total é simétrica: estados de spin
simétricos com $J$ par (orto, peso 6), antissimétricos com $J$ ímpar (para, peso
3). O orto-deutério, que contém $J = 0$, é estável a 0 K; à temperatura ambiente,
orto:para = 2:1.
\end{solution}

\begin{solution}{exo:b3:statistical-thermo-applied:5}
$\Delta_rE_0 = 2 \times 1883.7 - 2170.3 - 1542.3 = \qty{54.8}{cm^{-1}}$; $\eu^{-54.8/207.2} = 0.768$, e
$4 \times 0.768 = 3.07$. O valor completo, 3.26, é 6\,\% maior: os fatores translacional,
rotacional e vibracional só se cancelam exatamente no limite clássico, e a
rotação dessas moléculas leves ainda é parcialmente quantizada a \qty{298}{K}.
\end{solution}

\begin{solution}{exo:b3:statistical-thermo-applied:6}
\ce{N2O}: duas orientações, $R\ln2 = \qty{5.76}{J.K^{-1}.mol^{-1}}$. \ce{CH3D}: quatro, $R\ln4 =
\qty{11.53}{J.K^{-1}.mol^{-1}}$ (limites superiores, atingidos se os arranjos forem aleatórios).
\end{solution}

\begin{solution}{exo:b3:statistical-thermo-applied:7}
$x = 797.6/300 = 2.659$: função de Einstein $x^2\eu^{-x}/(1 - \eu^{-x})^2 = 0.572$; $C_V/R = 3.072$,
$C_p = 4.072R = \qty{33.86}{J.K^{-1}.mol^{-1}}$, 0.4\,\% abaixo da tabela (anarmonicidade).
\end{solution}

\begin{solution}{exo:b3:statistical-thermo-applied:8}
$x = 1.029$: $x^2\eu^{-x}/(1 - \eu^{-x})^2 = 0.916$, isto é, \qty{7.62}{J.K^{-1}.mol^{-1}}, 92\,\% de $R$: a
vibração do iodo é quase clássica à temperatura ambiente.
\end{solution}

\begin{solution}{exo:b3:statistical-thermo-applied:9}
B é mais rica em \ce{^{18}O}. $\alpha_{\mathrm{A/B}} = (1 - 0.0100)/(1 + 0.0020) = 0.98802$.
\end{solution}

\begin{solution}{exo:b3:statistical-thermo-applied:10}
$\Delta\mathrm{ZPE} = \frac12(1 - 1/\sqrt2)(2000 - 3000) = \qty{-146}{cm^{-1}}$; $K = \eu^{146.4/207.2} = 2.0$.
O deutério vai preferencialmente para X, a ligação mais rígida, onde a economia de
\omterm{def:b3:quantum-model-systems:oscillator}{energia do ponto zero} é maior.
\end{solution}

\begin{solution}{exo:b3:statistical-thermo-applied:11}
$\Delta_rH^\circ = R\ln10\,(-1.766 + 3.392)/(1/900 - 1/1100) = 19.145 \times 1.626/\num{2.020e-4} =
\qty{154}{kJ/mol}$. Os dois átomos carregam $2 \times \frac52RT$ de entalpia; a molécula, $\frac52RT + RT$
(rotação) mais sua energia vibracional $R\theta_{\mathrm V}/(\eu^{\theta_{\mathrm V}/T} - 1)$, um pouco menos que $RT$:
a diferença, cerca de \qty{5}{kJ/mol} a \qty{1000}{K}, soma-se a $D_0$.
\end{solution}

\begin{solution}{exo:b3:statistical-thermo-applied:12}
Níveis 0 ($g = 1$) e $6hcB$ ($g = 5$): $q = 1 + 5\eu^{-x}$, $\langle\varepsilon\rangle/kT = 5x\eu^{-x}/q$, $\langle\varepsilon^2\rangle/(kT)^2 =
5x^2\eu^{-x}/q$, e $C/R = \langle\varepsilon^2\rangle/(kT)^2 - (\langle\varepsilon\rangle/kT)^2 = 5x^2\eu^{-x}/(1 + 5\eu^{-x})^2$. A
\qty{100}{K}, $x = 5.121$: $C/R = 0.783/1.061 = 0.738$.
\end{solution}

\begin{solution}{pb:b3:statistical-thermo-applied:1}
\textbf{1.} $2 \times 2 = 4$: três simétricos (o tripleto), um antissimétrico (o singleto).
\textbf{2.} Os prótons são férmions, logo a função de onda total muda de sinal na
permutação; a função rotacional dá $(-1)^J$: os $J$ pares acompanham o singleto
antissimétrico (para), os $J$ ímpares o tripleto simétrico (orto).
\textbf{3.} Níveis pares 1, níveis ímpares 3.
\textbf{4.} A \qty{300}{K} ($T = 3.5\,\theta_{\mathrm R}$), as somas rotacionais par e ímpar são quase iguais;
ponderadas por 1 e 3, dão 1:3 (fração para 0.2507).
\textbf{5.} $F(1) = 2 \times 59.322 - 4 \times 0.0471 = \qty{118.46}{cm^{-1}} = \qty{1.417}{kJ/mol}$.
\textbf{6.} $1/(1 + 9\eu^{-170.4/20.37}) = 1/(1 + 9 \times \num{2.3e-4}) = 0.998$.
\textbf{7.} Com apenas $J = 0$ e 1, metade para significa $9\eu^{-2\theta_{\mathrm R}/T} = 1$, $T = 2\theta_{\mathrm R}/\ln9 =
0.91\,\theta_{\mathrm R} = \qty{78}{K}$: a competição se dá entre o peso 9 de $J = 1$ e seu
fator de Boltzmann.
\textbf{8.} 0.75. Sem catalisador, a conversão leva dias; a liquefação,
horas.
\textbf{9.} $(0.75 - 0.002) \times 1.417 = \qty{1.060}{kJ/mol}$.
\textbf{10.} $1.060/\num{2.016e-3} = \qty{526}{kJ/kg}$.
\textbf{11.} $0.905/\num{2.016e-3} = \qty{449}{kJ/kg}$.
\textbf{12.} O calor de conversão vale 1.17 vez a entalpia de vaporização.
\textbf{13.} Ele é produzido dentro do líquido, e não trazido pelas paredes.
\textbf{14.} $x_{\mathrm o}(t) = x_{\mathrm{eq}} + (0.75 - x_{\mathrm{eq}})\eu^{-kt}$, com $x_{\mathrm{eq}} = 0.002$.
\textbf{15.} $0.002 + 0.748\eu^{-0.274} = 0.571$.
\textbf{16.} $(0.75 - 0.571) \times 1.417 = \qty{0.254}{kJ/mol}$.
\textbf{17.} $0.254/0.905 = 0.28$: 28\,\% do tanque em um dia.
\textbf{18.} Depois de \qty{168}{h}, $x_{\mathrm o} = 0.112$, calor de \qty{0.904}{kJ/mol}, igual à entalpia de
vaporização: nesse modelo, o tanque fica vazio. A estimativa exagera a perda
porque as moléculas evaporadas levam consigo seu \omterm{def:b3:statistical-thermo-applied:spin-isomers}{orto-hidrogênio} não
convertido; a perda real é menor, mas ainda grande.
\textbf{19.} $t_{1/2} = \ln2/0.0114 = \qty{61}{h}$.
\textbf{20.} No liquefator, sobre leitos de catalisador em temperaturas sucessivas, para
que o refrigerador remova o calor de conversão antes de o líquido ser armazenado.
\textbf{21.} Sim: a fração para de equilíbrio a \qty{300}{K} é 0.2507.
\textbf{22.} No hidrogênio normal, as duas formas não podem se interconverter, de modo
que a rotação de cada uma congela em sua própria escada de níveis; no hidrogênio
de equilíbrio, parte do calor fornecido converte para em orto.
\textbf{23.} Cerca de 3.57, perto de \qty{50}{K}: além de excitar a rotação, o calor converte
moléculas de $J = 0$ em $J = 1$, \qty{118}{cm^{-1}} acima, uma reação de entalpia
positiva.
\textbf{24.} 1.50: a rotação está congelada nas duas formas (para em $J = 0$, orto em $J = 1$).
\textbf{25.} \textbf{Calor de conversão / entalpia de vaporização $\approx 1.2$ para o hidrogênio
normal}: só sua conversão poderia evaporar todo o tanque, e é por isso que o
hidrogênio é convertido em para antes do armazenamento.
\end{solution}
